\begin{theorem}[Nonvanishing of local constants]\label{mod:delta-nonvanishing}
Let $F$ be a nonarchimedean local field, let $\chi$ and $\psi$ be exact
multiplicative and additive character data, and let $\gamma$ be an admissible
denominator.  By Definition~\ref{mod:delta-finitedefinition},
\[
 \Delta_F^{\mathrm{fin}}(\chi,\psi;\gamma)
 =\chi(\gamma)\,
   \operatorname{ph}\bigl(G_\gamma(\chi,\psi)\bigr).
\]
The character value $\chi(\gamma)$ is nonzero because the quasi-character is
$\mathbf C^\times$-valued, and the total phase is nonzero at every complex
argument by Definition~\ref{mod:basic-phase}.  Hence
\[
 \Delta_F^{\mathrm{fin}}(\chi,\psi;\gamma)\ne0.
\]
Moreover, the theorem
\texttt{finiteGaussSum\_ne\_zero} proved in
Definition~\ref{mod:delta-finitedefinition} shows that
$G_\gamma(\chi,\psi)$ itself is genuinely nonzero.  Thus the zero branch of
the totalized phase is not used, although only the unconditional
nonvanishing of the total phase is needed for the displayed product.
\LeanFile{LanglandsFirstMainLemma/Delta/Nonvanishing.lean}
\lean{LanglandsFirstMainLemma.deltaFinite_ne_zero}
\leanok
\SourceText{Definition~\ref{eq:def-delta}; Corollary~\ref{cor:local-gauss-integral-nonzero}.}
\uses{mod:delta-finitedefinition}
\FormalizationContract The exported theorem has no conductor-case hypotheses
and assumes no nonvanishing fact: it combines the nonzero character value
with the unconditional nonvanishing of the total phase.  The separately
proved nonvanishing of the underlying finite Gauss sum records that the
phase's zero branch is never reached.
\end{theorem}
